% Generated draft: requires mathematical and source review.
\documentclass[11pt]{article}
\usepackage[margin=25mm]{geometry}
\usepackage{fontspec}
\setmainfont{Latin Modern Roman}
\usepackage{amsmath,amssymb,mathtools,unicode-math}
\setmathfont{Latin Modern Math}
\usepackage{enumitem,needspace,xurl,hyperref}
\setlength{\parindent}{0pt}
\setlength{\parskip}{5pt}
\setlength{\emergencystretch}{3em}
\newcommand{\N}{\mathbb{N}}
\newcommand{\Z}{\mathbb{Z}}
\newcommand{\Q}{\mathbb{Q}}
\newcommand{\R}{\mathbb{R}}
\newcommand{\C}{\mathbb{C}}
\newcommand{\F}{\mathbb{F}}
\newcommand{\A}{\mathbb{A}}
\newcommand{\PP}{\mathbb{P}}
\newcommand{\E}{\mathbb{E}}
\newcommand{\eps}{\varepsilon}
\newcommand{\dd}{\,\mathrm{d}}
\begin{document}
\section*{UnsolvedMath ID: 30005131}
\textbf{Source problem code:} catalogue entry
\subsection{Definitions and mathematical statement}
Let \(X/\Z\) be a proper regular arithmetic scheme and let \(n\in\Z\). The record does not specify the precise convention for an arithmetic scheme or all hypotheses required below; the conjecture must therefore be read subject to the unresolved qualifications stated here.

Write \(\zeta(X;s)\) for the zeta function of \(X\). Under the usual finite-type interpretation, its defining Euler product in a region of convergence is
\[
\zeta(X;s)=\prod_{x\in |X|}\bigl(1-\#\kappa(x)^{-s}\bigr)^{-1},
\]
where \(|X|\) denotes the set of closed points and \(\kappa(x)\) their finite residue fields. To define a leading value at \(n\), assume that this function has a nonzero meromorphic continuation to a neighborhood of \(n\). If
\[
\zeta(X;s)=(s-n)^{r_{X,n}}\bigl(a_{X,n}+O(s-n)\bigr),
\qquad r_{X,n}\in\Z,\quad a_{X,n}\ne0,
\]
set \(\zeta^*(X;n)=a_{X,n}\). This uses a Laurent expansion when \(n\) is a pole. The original excerpt says ``Taylor development'' and does not specify whether poles are excluded or a Laurent interpretation is intended.

The excerpt invokes, under unspecified assumptions, a perfect complex of abelian groups
\[
K_{X,n}=R\Gamma_{W,c}(X;\Z(n)),
\]
and a perfect complex
\[
D_{X,n}=R\Gamma\bigl(X;L\Omega^{<n}_{X/\mathbb S}\bigr).
\]
Here \(\mathbb S\) denotes the sphere spectrum, and the excerpt describes the latter complex as derived de Rham cohomology relative to the sphere spectrum, modulo the \(n\)-th step of the Hodge filtration. Its exact construction, including its convention for arbitrary negative \(n\), is not supplied. The source attributes the construction to the use of work of Antieau and Bhatt--Morrow--Scholze.

For a perfect complex \(P\) over \(\Z\), its determinant line is
\[
\det_{\Z}P=\bigotimes_i\left(\bigwedge^{\operatorname{rk}P^i}P^i\right)^{(-1)^i},
\]
using a bounded finite projective representative; exponent \(-1\) denotes the dual line. The displayed target in the extracted source appears to give the integral line
\[
L_{X,n}=\det_{\Z}K_{X,n}\otimes_{\Z}\det_{\Z}D_{X,n}
\]
and, when Beilinson's regulator behaves as expected, a canonical real isomorphism
\[
\lambda_{X;n}:\R\xrightarrow{\sim}L_{X,n}\otimes_{\Z}\R.
\]
The regulator condition and the hypotheses ensuring these constructions are not stated. The extracted display is damaged, so its tensor and dual conventions require confirmation.

\textbf{Conjectural target, subject to these qualifications.} For every such \(X\) and every \(n\in\Z\) for which the requisite complexes, regulator isomorphism, and leading value are defined under the intended source hypotheses, the supplied record asks whether
\[
\det(\lambda_{X;n})=\pm\zeta^*(X;n).
\]
This is the source's symbolic formulation. To make its scalar meaning explicit, a generator \(\ell\) of \(L_{X,n}\) determines a scalar \(a\) by \(\lambda_{X;n}(1)=a\ell\), well defined up to sign. If the source uses this scalar convention, the displayed conjecture means \(a=\pm\zeta^*(X;n)\). The record does not establish whether this convention or the reciprocal convention is intended; the normalization must be checked before treating that interpretation as a precise formulation.

The original excerpt also contains a damaged assertion labeled ``Theorem'' involving \(\varepsilon_{X;n}\), archimedean leading values, \(A(X)\), and \(d\). These quantities and their relationships to the conjecture are not sufficiently defined to reconstruct that assertion. It is retained as an unresolved source qualification, not imposed as an additional hypothesis or asserted here.
\subsection{Short English statement}
The supplied conjecture relates a regulator-induced determinant trivialization for a proper regular arithmetic scheme to the leading coefficient of its zeta function at an integer. Its exact hypotheses, determinant normalization, and some extracted formulas remain unresolved. The supplied status is partially solved; this is a statement-only draft requiring mathematical review.
\subsection{Sources}
\begin{itemize}
\item Local registry: \texttt{data/unsolvedmath/problems.json}, with \texttt{catalog/problems.json} as the supplementary registry; supplied permanent UnsolvedMath ID \(30005131\), exact source problem code \texttt{catalogue entry}, title ``Determinant Formula for Leading Zeta Values.'' No individual proposer is supplied; the proposed year is recorded as \(2022\).
\item Supplied report citation: \emph{Algebraic K-Theory (2023)}, Oberwolfach Reports, DOI \texttt{10.4171/owr/2022/24}, pp.\ 1313--1376. \url{https://doi.org/10.4171/owr/2022/24}. The supplied title year \(2023\), DOI containing \(2022\), and proposed year \(2022\) are retained without reconciliation.
\item Supplied literature source labeled ``Recent de Rham--syntomic paper'': \url{https://arxiv.org/abs/2603.21471}. The record supplies no authors or full title for this reference.
\item Supplied literature source labeled ``Oberwolfach report'': \url{https://doi.org/10.4171/owr/2022/24}.
\item Literature assessment dated exactly \texttt{2026-08-22}: the supplied curation labels the problem partially solved and reports special cases and compatibility results, a recent computation of the de Rham correction factor, and an unproved full formula for arbitrary proper regular arithmetic schemes. This is supplied dated curation, not independently verified literature assessment; the background expressly requests independent verification of time-sensitive status and cited claims.
\item The registry labels the statement \texttt{corrected\_verified} and reports that the official report was checked directly. This is a supplied verification claim, not verification performed for this draft. The excerpt attributes underlying work to Antieau and Bhatt--Morrow--Scholze without providing separate bibliographic references.
\end{itemize}
\end{document}
